\documentclass[11pt,a4paper]{article}
\usepackage[T1]{fontenc}
\usepackage{lmodern,amsmath,amssymb,amsthm,microtype,booktabs}
\usepackage[margin=25mm]{geometry}
\usepackage[colorlinks=true,urlcolor=blue,linkcolor=blue]{hyperref}
\newtheorem{proposition}{Proposition}
\newcommand{\B}{\mathfrak B}
\newcommand{\dd}{\mathop{}\!\mathrm d}
\title{Removable energy pinches in the recoil graph\\and a direct two-point calibration}
\author{Alan Day \and Chengheng\thanks{AI research assistant (OpenAI).}}
\date{Working research note --- 31 October 2026}
\begin{document}\maketitle
\begin{abstract}
The internal Barnes integral of the pure-recoil tetrahedron-complement
graph has energy-transfer pinch poles, including fourth-order poles
on its symmetric slice. We regularize all twelve left--right Gamma
pairs and prove that the resulting function is entire in the energy
transfer. The complete vertex product cancels the pinch factors
exactly after Schwarz pairing. A separate fixed-mass derivation of
the direct \(C_0\) two-point factor retains energy poles.
These are statements about specified recoil-limit functions, not
a physical four-point insertion prescription or a finite-\(G\) result.
\end{abstract}

\section{Object and pinch divisors}
We use the CLP graph \cite{clp}, at radius one, taking
\(\kappa\to\infty\) at fixed real masses first. Complex masses below
mean continuation of that limiting function. No interchange with a
finite-\(\kappa\) growing-mass limit is asserted.
Set
\begin{align}
\boldsymbol\mu&=(m_1,\nu,m_2,m_3,\nu,m_4),&
m_{1,3}&=M+\omega/2\pm u,&
m_{2,4}&=M-\omega/2\pm v .
\end{align}
The vertex sets are \(123,156,246,345\), and the quadrilateral sets
are \(1245,1346,2356\). Their half-sums are \(a_j\) and \(r_l\);
\(\sum_j a_j=\sum_l r_l\). The previously derived recoil integral is
\begin{equation}
\B=\frac1{2\pi i}\int_{d-i\infty}^{d+i\infty}
e^{-i\pi z}
\frac{\prod_{j=1}^4\Gamma(z+ia_j)
\prod_{l=1}^3\Gamma(\tfrac12-z-ir_l)}
{\Gamma(z+\tfrac12)}\,\dd z,\qquad 0<d<\tfrac12.
\label{eq:barnes}
\end{equation}
This line applies to the real starting data. Thereafter the contour
must separate the transported pole families.
A left pole \(-ia_j-n\) and a right pole
\(\tfrac12-ir_l+m\) meet when
\begin{equation}
p_{jl}:=\tfrac12+i(a_j-r_l)=-N,\qquad N=n+m\geq0.
\label{eq:pinch}
\end{equation}
For equal probes,
\[
a_j=M+\nu/2+(\pm u\pm v)/2,\quad
r_s=M+\nu+s\omega/2\ (s=\pm1),\quad r_0=2M.
\]
Writing \(w_j=M+\nu-a_j\), the moving divisors are
\begin{equation}
p_{js}=\tfrac12-iw_j-is\omega/2,\qquad
\omega_{jsN}=-s[\,2w_j+i(2N+1)\,].
\label{eq:locations}
\end{equation}
The signs are correlated. For \(w_j>0\), the positive-real-part
upper-half-plane poles belong to \(s=-1\); the opposite chirality
has the conjugate families. The \(r_0\) divisors are fixed and
nonsingular for real parameters. If \(q\) of the \(a_j\)'s coincide,
the corresponding pole order is at most \(q\), and is \(q\) when
its leading coefficient does not vanish.

\section{An entire regularized function}
Define
\begin{equation}
G=\prod_{j,l}\Gamma(p_{jl}),\qquad H=\B/G .
\end{equation}
For distinct real \(a_j\), closing \eqref{eq:barnes} to the left gives
\begin{align}
H={}&\sum_j A_j
\prod_{\substack{k\ne j\\l=1,2,3}}\frac1{\Gamma(p_{kl})}\,
{}_4F_3\left(\begin{matrix}
\tfrac12+ia_j,p_{j1},p_{j2},p_{j3}\\
(1+i(a_j-a_k))_{k\ne j}
\end{matrix};1\right),\label{eq:left}\\
A_j={}&
\frac{e^{-\pi a_j}\prod_{k\ne j}\Gamma(i(a_k-a_j))}
{\Gamma(\tfrac12-ia_j)} .
\end{align}
Parenthesized subscript lists in the hypergeometric denominator
denote the listed parameters, not shifted factorials.
The excess (sum of lower parameters minus sum of upper parameters)
is exactly one in every term.

\begin{proposition}
For fixed real \(M,\nu,u,v\), the continuation of \(H\) from real
\(\omega\) is entire in \(\omega\), including coincident \(a_j\)'s.
\end{proposition}
\begin{proof}
For distinct real \(a_j\), the lower parameters in \eqref{eq:left}
have real part one and are independent of \(\omega\).
The two varying upper parameters have constant sum.
Gamma-ratio estimates give a uniform \(O(n^{-2})\) bound for the
series terms on every compact \(\omega\) set. Finite terms are
polynomials in the upper parameters, including their terminating
values. Thus the series is entire in \(\omega\), as are the reciprocal
Gamma factors. This uses ordinary hypergeometric convergence and
Gamma properties \cite{dlmf16,dlmf5}.

For coincident \(a_j\), perturb real \(u,v\) in a generic direction.
Near zero perturbation the expression has only finite-order
Laurent singularities: the lower hypergeometric parameters remain
near one, while only the Gamma differences can have poles.
Each Laurent coefficient is entire in \(\omega\).
Near real \(\omega\), the original vertical integral and \(G^{-1}\)
are holomorphic in the perturbation, with a common exponential
tail bound. Hence all negative Laurent coefficients vanish there.
The identity theorem makes them vanish for every \(\omega\).
The coincident limit is therefore removable.
\end{proof}

For computation near coincident \(a_j\), right residues provide
a better representation:
\begin{align}
H={}&\sum_l D_l
\prod_{\substack{j=1,2,3,4\\h\ne l}}\frac1{\Gamma(p_{jh})}\,
{}_4F_3\left(\begin{matrix}p_{1l},p_{2l},p_{3l},p_{4l}\\
1-ir_l,(1+i(r_h-r_l))_{h\ne l}
\end{matrix};1\right),\label{eq:right}\\
D_l={}&\frac{-i e^{-\pi r_l}
\prod_{h\ne l}\Gamma(i(r_l-r_h))}{\Gamma(1-ir_l)} .
\end{align}
This formula is used away from collisions of the right families;
one takes the full continued sum at any apparent singularity.
The \(-i\) includes \(e^{-i\pi/2}\); the clockwise closing sign
cancels the sign of the right Gamma residue.
Reflection of the three right Gamma factors leaves a \(z^{-2}\)
Gamma ratio. On closing lines chosen away from poles the sine
factors give exponential decay in both imaginary directions.
The remote line therefore vanishes. The excess is again one.
At a pinch, one or more upper parameters are nonpositive integers,
making the relevant term terminate.

For a generic simple pole, \(\partial_\omega p_{js}=-is/2\).
Its residue is explicitly
\begin{equation}
\mathop{\rm Res}_{\omega_{jsN}}\B
=\frac{2is(-1)^N}{N!}\,
 A_j\!\prod_{h\ne s}\Gamma(p_{jh})\,
{}_4F_3\left(\begin{matrix}\tfrac12+ia_j,(p_{jl})_{l=1,2,3}\\
(1+i(a_j-a_k))_{k\ne j}\end{matrix};1\right)_{\!p_{js}=-N}.
\label{eq:residue}
\end{equation}
There is no division by a vanishing leading value in our general
entireness argument.

\section{Exact cancellation in the complete graph}
The continued product of the four recoil vertices is built from
\begin{equation}
T(\mu_1,\mu_2,\mu_3)=
\pi^{-3/2}\sqrt{\prod_{e=1}^3\frac{\mu_e}{\sinh\pi\mu_e}}\,
\prod_{e=1}^3\cosh\frac{\pi(S-2\mu_e)}2,\qquad S=\sum_e\mu_e.
\end{equation}
The twelve quantities \(r_l-a_j\) are precisely the twelve
half triangle deficits \((S_v-2\mu_e)/2\).
Each edge occurs twice, so the square roots disappear from the
full vertex product. With
\begin{equation}
L(\boldsymbol\mu)=\prod_{e=1}^6\frac{\mu_e}{\sinh\pi\mu_e},
\end{equation}
we obtain
\begin{equation}
\prod_v T_v=\pi^{-6}L\prod_{j,l}\sin(\pi p_{jl}).
\label{eq:vertex}
\end{equation}
Schwarz continuation is \(f^\sharp(z)=\overline{f(\bar z)}\);
at fixed real transverse data \(p^\sharp=1-p\).
The reflection formula then gives the exact identity
\begin{equation}
\boxed{\quad
\mathcal Z_{\rm rec}:=(\prod_v T_v)\B\B^\sharp
=\pi^6 L\,H H^\sharp .
\quad}\label{eq:regulargraph}
\end{equation}
Indeed \(GG^\sharp=\pi^{12}/\prod\sin(\pi p_{jl})\).
All the internal pinch Gamma factors cancel, with their full
multiplicities. A normalized graph is obtained by dividing both
sides by the same nonzero reference value. This does not introduce
a physical disconnected normalization.

The regular value at a former pinch need not vanish.
Cancellation of poles also does not remove either large-mass
Bessel branch found in the earlier paired-graph analysis.
It reorganizes the same function, rather than changing it.

For fixed positive \(\nu\), the only remaining \emph{possible}
energy poles are those of the four observer-edge factors:
\begin{equation}
\omega=-2(M\pm u)+2in,\qquad
\omega=2(M\pm v)-2in,\qquad n\in\mathbb Z\setminus\{0\}.
\end{equation}
The ratio \(\mu/\sinh\pi\mu\) is regular at \(\mu=0\).
Some of the displayed poles may be removed by zeros of \(H H^\sharp\);
we do not classify them here.
In particular the graph is holomorphic throughout
\begin{equation}
|\mathop{\rm Re}\omega|<2[\,M-\max(|u|,|v|)\,],
\qquad M>\max(|u|,|v|).
\label{eq:strip}
\end{equation}
If \(M>2\nu\) and \(|u|+|v|<\nu\), all internal pinch columns
lie in this strip. Holomorphy alone does not equate infinite
integration cycles: their ends and the external weights still
require control.

\section{A fixed-mass direct-\(C_0\) calibration}
CLP (3.35) uses a direct \(C_0 C_0^\sharp\) factor, observer
weights and an insertion normalization \(\mathcal N\) \cite{clp}.
Its fixed-angle heavy-heavy-light prescription is not uniform
in the fixed-mass recoil limit. We instead derive the
mass-dependent ratio directly from the same small-\(b\) limit.

For a triple of masses write
\[
\bar z_0=(1-iS)/2,\quad
\bar z_a=(1+i(S-2\mu_a))/2,\quad z_a=1-\bar z_a,\quad L_\kappa=-i\kappa .
\]
The four double-sine arguments in \(V^2\) are
\(b\bar z_0-1/b\) and \(b\bar z_a+1/b\).
Using the exact shift relation and its fixed-parameter limit,
\[
S_b(bz+1/b)=2\sin(\pi z)S_b(bz),\qquad
[2\sin\pi(z-L_\kappa)]^{-1}\sim-i e^{-\pi\kappa+i\pi z},
\]
gives the mass-dependent recoil factor
\begin{equation}
v(\boldsymbol\mu)=
e^{\pi S/2}\Gamma(\bar z_0)\prod_a\Gamma(z_a)^{-1}.
\end{equation}
Here \(\bar z_0+\sum_a\bar z_a=2\), so the other small-\(b\)
powers cancel in reference-mass ratios.
The previously established \(U=V^2/C_0\) factor tends, in the
same ratio sense, to
\[
F(\boldsymbol\mu)=
\frac{\prod_a\Gamma(1-i\mu_a)}{\prod_a\Gamma(z_a)^2}.
\]
Consequently
\begin{align}
c(\boldsymbol\mu)&=
e^{\pi S/2}
\frac{\Gamma(\tfrac12-\tfrac{iS}{2})
 \prod_a\Gamma(\tfrac12-\tfrac i2(S-2\mu_a))}
{\prod_a\Gamma(1-i\mu_a)},\label{eq:c}\\
D(\boldsymbol\mu):=c c^\sharp&=
\frac{\pi e^{\pi S}\prod_a[\sinh(\pi\mu_a)/\mu_a]}
{\cosh(\pi S/2)\prod_a\cosh(\pi(S-2\mu_a)/2)}.\label{eq:d}
\end{align}
Only reference ratios of \(v,c,D\) are fixed by the limiting
derivation; common mass-independent factors have been removed.
This is not an absolute normalization of physical insertions.

For \(\mu_{1,2}=M\pm\delta/2,\ \mu_3=\nu,\ M>2\nu>0\),
the direct factor \(D\) has simple poles at
\(\delta=\pm\nu+i(2N+1)\), \(N\in\mathbb Z\).
At \(\delta_0=\nu+i(2N+1)\) its nonzero residue is
\begin{equation}
\mathop{\rm Res}_{\delta_0}D
=-\frac{2e^{\pi(2M+\nu)}
 \sinh(\pi\mu_1)\sinh(\pi\mu_2)}
{\mu_1\mu_2\nu\cosh(\pi(2M+\nu)/2)\cosh(\pi(2M-\nu)/2)}
\bigg|_{\delta=\delta_0}.
\label{eq:cres}
\end{equation}
Here \(\delta\) is a two-point mass difference and must be mapped
explicitly before comparison with a four-point global transfer.

This contrast supplies a calibration condition, not a unique
sewing rule. The complex continuation of \(\mathcal N\), the
physical inner product, the spectral measure and disconnected
factorization remain to be specified. One cannot restore desired
poles by arbitrarily multiplying a graph by \(D\) or its square root.
The mixed representation and contour issues of CSTY \cite{csty}
therefore remain relevant.

\section{Numerical evidence and continuation}
The saved driver performs 74 checks at 40 decimal digits, with
selected 60-digit comparisons. It compares the two residue
representations, the unregularized four-term sum, an original
vertical integral, and the complete vertex identity.
Three examples have left-sum condition numbers up to
\(3.68\times10^4\), while the corresponding right sums have
condition numbers at most \(1.0675\).

For \(M=3,\nu=1\), the transverse pairs
\((1/8,1/4),(0,1/4),(0,0)\) give simple, double and quadruple
pinches. Both \(N=0,1\) are tested.
Approaches \(\omega=\omega_0+\epsilon\), with
\(\epsilon=10^{-3},10^{-5},10^{-7}\), verify the predicted
\(\epsilon^{-q}\) behavior of \(\B\) and a finite full graph.
The largest final relative approach errors are approximately
\(8.12\times10^{-7}\) for the pole coefficient and
\(5.31\times10^{-7}\) for the full graph.
These are finite-\(\epsilon\) effects, not quadrature error bars.
For example, at the symmetric fourth-order pinch \(\omega_0=1+i\),
\[
\mathcal Z_{\rm rec}(\omega_0)=
-4.38134684304\times10^{-10}
+1.73872354359\times10^{-10}i .
\]
The exact-pinch value is evaluated using the terminating
regularized representation, not inferred solely by extrapolation.

Two mass samples also test the exact double-sine \(V^2\)
factor at \(\kappa=40,160,640\) against \(v\).
This tests a constituent of the \(C_0\) derivation, not the
complete finite-\(\kappa\) \(C_0\) or four-point graph.
No interval certificate or finite-\(G\) dynamical claim is made.

The next task is to calibrate observer weights and disconnected
limits against \eqref{eq:d}, then seek an actual sewing/inner-product
prescription. Identification with a known Racah function requires
a parameter, measure and normalization dictionary, not just a
similar Barnes integral. Internal pinch locations alone no longer
select a physical external cycle; growth at infinity and physical
insertions must do that work.

\paragraph{Reproducible records.}
Chinese derivation: notes/stage4-recoil-poles.md.
Implementation: code/recoil\_poles.py and code/validate\_recoil\_poles.py.
Raw values, precision, condition numbers and input hashes:
results/stage4/recoil-poles/diagnostics.json.

\begin{thebibliography}{9}
\bibitem{clp}
J.~Chandra, D.~Liska and B.~Post,
\emph{3D de Sitter from complex Virasoro TQFT: Observers and knots},
\href{https://arxiv.org/abs/2610.30247v1}{arXiv:2610.30247v1},
especially (3.17), (3.35)--(3.37), (3.43), (3.46)--(3.47) and Appendix A.
\bibitem{csty}
Y.~Chen, D.~Stanford, H.~Tang and Z.~Yang,
\emph{Negative shocks versus static patch holography},
\href{https://arxiv.org/abs/2607.14042v1}{arXiv:2607.14042v1}.
\bibitem{dlmf5}
NIST Digital Library of Mathematical Functions,
\href{https://dlmf.nist.gov/5.2}{\S5.2} and
\href{https://dlmf.nist.gov/5.5}{\S5.5}.
\bibitem{dlmf16}
NIST Digital Library of Mathematical Functions,
\href{https://dlmf.nist.gov/16.2}{\S16.2} and
\href{https://dlmf.nist.gov/16.5}{\S16.5}.
\end{thebibliography}
\end{document}
